\scanpage{255}

\mspage{130}

qui \inserted{$(\unsure{c})$} est surjectif ssi on est dans la ``situation d'isomorphie'' i.e. $\delta:G\unsure{\longrightarrow|\longleftrightarrow} Z$ surjectif.

% margin note written sideways
\margin{\uline{NB} On a des iso canoniques $(115')\left\{\begin{aligned}T_0&\xrightarrow[\sim]{\unsure{\kappa_0|\pi_0}}\widetilde Z\\T_1&\xrightarrow[\sim]{\unsure{\kappa_1|\pi_1}}\widetilde Z\end{aligned}\right.$ mais on se gardera d'identifier $T_0,T_1$ (sous-groupes distincts de $G$) à $\widetilde Z$ (groupe quotient de $G$)}

On a un diagramme symétrique à (113) en y remplaçant l'indice $0$ par $1$, il est déduit de (113) en \deleted{y \unsure{remplaçant|transformé}} \inserted{le transformant} par l'autom $\operatorname{int}(\sigma)$, où $\sigma$ est considéré comme élément de $G$.

\deleted{Si on veut une description ``symétrique'' de $G$, \unsure{on} on peut écrire}

\[G=\deleted{\unsure{B}}T_G.H,\quad\text{avec }T_G\cap H=T_H\quad(98)\tag*{(116)}\]

donc dans cette description \inserted{de $G$}, il s'agit de bien expliciter

\[\left\{\begin{aligned}
&\text{a) le groupe }T_G,\quad\text{b) le plongement de }T_H\text{ dans }T_G,\text{ et}\\
&\text{\deleted{enfin} c) l'action de }T_G\text{ sur }H,\quad\text{\deleted{enfin} ce qui}\\
&\text{permet de reconstituer l'inclusion}\\
&\hspace{8em}H\subset G\\
&\text{de (64), il faut de plus}\\
&\text{d) expliciter l'hom d'inclusion } Z\subset T_G,\text{ et enfin}\\
&\text{e) expliciter l'hom }\delta|T_G:T_G\longrightarrow Z,\text{ pour}\\
&\text{réconstituer le diagramme (64).}
\end{aligned}\right.\tag*{(116')}\]

Nous allons expliciter $T_G$ \deleted{\unsure{comme} à l'aide de la suite exacte} comme coextension centrale

\[\deleted{1\longrightarrow\unsure{\mathfrak{z}|\mathfrak{Z}}\longrightarrow
\underbrace{T_{0\deleted{\unsure{G}}}\times T_1}_{\simeq\,\widetilde Z\times\widetilde Z}
\longrightarrow T_G\longrightarrow 1,}\tag*{\deleted{(117)}}\]

\deleted{en supposant pour simplifier que \unsure{$\widetilde Z_0$|$\widetilde T_0$} est contenu \inserted{(donc $T_0,T_1$)} dans $T_1$ aussi commutatifs, donc aussi \unsure{$Z$}} \inserted{\deleted{$T_0$ et $T_1$ isom à $\widetilde Z$ sont}}

\[\deleted{T_G\simeq T_H\times T_0\simeq Z\times\widetilde Z}\]

% archive sheet number: 254
